\chapter[不要忘记]{不要忘记}

\chaptersubtitle{《魔法少女和美》}

在《魔法少女小圆》电视动画的发展过程中，确立了三条主要原则；这三条原则合在一起，对本部分将要讨论的三部外传作品都构成了某种挑战。
第一，从魔法少女的起源开始，魔法少女就是，而且一直都是魔女。
考虑到魔女原本是父权制对女性力量恐惧的形象\cite{ref103}，这意味着《魔法少女小圆》中的魔法少女既内在地挑战着支配性的权力结构，也很容易被解读为某种腐败且会使他人腐败的因素。
第二，《魔法少女小圆》的世界由孵化者的系统所支配；在这个系统中，宇宙的熵通过一种未经魔法少女同意的方式被“清理”掉，也就是被转化为魔法少女堕落为魔女时的情感衰败。
换言之，支配性的权力结构剥削女性，并且对女性极其不公。
第三，这种权力结构高度稳定，直到小圆许下愿望、创造出她的新系统之前，它都没有被推翻。

暂且搁置小圆的系统是否比丘比的系统更加公正这一问题（但可参见间幕二和第十七章），第三条原则意味着，任何设定在动画主时间线中，或设定在丘比系统占据支配地位的任何“先前”时间线中的外传作品，最终都必须把丘比的系统视为一种无法逃脱、无法改变的事实。
再加上丘比的系统如何限制并剥削魔法少女的力量（第二条原则），以及这种力量又如何被等同于一种将女性力量描绘为腐败且危险的传统（第一条原则），外传作品便很容易落入这样一个陷阱：
把魔法少女描绘成危险的存在，而把丘比的系统描绘成用来遏制她们的必要之恶，甚至更糟，将其描绘为总体上的积极事物。

漫画《魔法少女和美》（平松正树编剧，天杉贵志作画）落入的，正是这个陷阱。

《魔法少女和美》的开篇中，标题角色和美几乎是一张彻底的白纸：
她赤身裸体，被封在一只箱子里，并且完全不记得自己的过去。
这是臭名昭著的 \emph{Turkey City Lexicon}\cite{ref104} 所称“白房间综合征”的一种极端形式：
故事开篇时，一个孤零零的角色身处白色房间，思考自己为何会来到这里；
这本词典相当不客气、但或许准确地假定，这是因为作者尚未想清楚故事的场景或角色。
这一批评是否适用于《魔法少女和美》并不明确。
一方面，和美的失忆承担着重要的情节功能，因为它使她能够成为观众发现她自身身份以及昴宿星圣团真相的媒介。
然而，《魔法少女和美》是一部连载时间相当长的作品，共有二十三话月刊连载，而且在电视动画完结之前很早便已开始。
《魔法少女和美》后期的许多揭示，在逻辑上都依赖于《魔法少女小圆》后期的揭示；
目前并不清楚，这些揭示之所以被推迟，是为了避免外传漫画剧透本篇动画的真相，还是因为平松和天杉并不知道相关剧透，不得不在连载中途改变漫画走向，以避免与本篇发生矛盾。

从这里开始，漫画揭示出和美与她的朋友们都是魔法少女，是一个七人团队“昴宿星圣团”的成员，并由一个看起来更友善的丘比变体“十兵卫”提供支持。
随着漫画展开，读者逐渐得知，昴宿星圣团远比最初看起来黑暗。
最终，故事揭示出她们原本是一个协作型魔法少女团体，而和美曾是她们的导师，但她最终变成了魔女。
她们合力将和美重新变成近似人类的形态，同时杀死并改造了当地的丘比个体，将其变成十兵卫；
与此同时，她们又用一道法术覆盖整座城市，使后来替补出现的丘比无法被察觉。
昴宿星圣团误以为十兵卫能够替她们净化灵魂宝石，但事实上，它只是制造出净化的幻觉；
在表象之下，她们仍然在缓慢地走向魔女化。
与此同时，和美只是复活原初和美的一连串尝试中最新的一次；
每一个复活体最终都会崩坏，并变成魔女。
在结局中，昴宿星圣团除两人以外全部死亡，并创造出一只巨大的复合魔女；
但和美开始意识到，她是一个不同于原初和美的人，于是与丘比签订契约，成为魔法少女，拯救城市和剩余的昴宿星圣团成员。

这个相当曲折的故事（由于背景故事并不是按事件发生的顺序讲述，而是按和美发现真相的顺序讲述，因而变得更加复杂）所做的，是将我在第九章讨论过的魔女与魔法少女的一体性字面化。
然而，这样做并不是为了尝试救赎魔女；
相反，它把魔法少女描绘成一种魔女。
在视觉上，这体现为和美的“日常”服装与她的魔法少女服装之间的对比。
前者在第二话中被描述为适合更年幼孩子的装束；
后者则比电视动画中的服装暴露得多，并且像是万圣节魔女服装的性化版本：
宽檐高顶的黑帽子、极短的黑裙、裸露的腹部、紧身上衣。
将其与小圆带有荷叶边的粉色裙装、沙耶香带有男孩气的铠甲，甚至电视动画中最性化的服装，也就是麻美高领束腰的连衣裙与打底裤，以及杏子大半敞开的卫衣和裤外短裙相比，漫画在天真孩童与强大、危险、具有性意味的魔法少女之间划出的区别便清晰可见。

这正是漫画与动画之间最大的差异：
在《魔法少女小圆》中，魔法少女与魔女的一体性，是关于希望与绝望循环的陈述的一部分；
而在《魔法少女和美》中，它却变成了对圣母/荡妇情结的拥抱。
在这一系列中，魔法少女全都隐藏着可怕的秘密：
要么她们本身就是魔女（和美），要么她们过去曾经做过可怕的事（昴宿星圣团）。
毫无例外，昴宿星圣团的每一个人，都被成为魔法少女之前的某种过错纠缠；
她们都参与欺骗和美，也都在其他魔法少女变成魔女之前协助追捕并囚禁她们。
这种转变并没有被描绘成我们在《魔法少女小圆》中看到沙耶香经历的那种逐渐沉入绝望的过程，而只是被描绘成使用过多魔法后会发生的事情；
换言之，魔女并不是绝望的产物，而是女性力量行使的产物。
更露骨的是，早期反派是一些拥有直接将女性转化为伪魔女之力量的魔法少女。

等到我们终于看到真正的魔女时，她们也与《魔法少女小圆》中的魔女相当不同。
在《魔法少女小圆》中，魔女被暗示为无心智、陷入妄想，并被困在自己的结界之中；
而在《魔法少女和美》中，她们能够直接与世界互动并体验世界，其中一些还拥有个性与智力。
因此，她们并不是浮士德式引用和女性绝望的具象化，而是魔法少女身体发生的怪诞变形。
本质上，在《魔法少女和美》中，魔法少女与魔女之间唯一的区别在于：
魔法少女是“可爱”或“性感”的，而魔女则是畸形的，或在身体上呈现为怪物。
换言之，魔女的丑陋并不来自她与丑陋情绪的关联，而在于她不再对男性凝视具有吸引力；
而这部漫画运用男性凝视的频率，远高于电视动画。

结局同样否定了《魔法少女小圆》中所描绘的魔女本质。
一旦和美与她幸存的两位朋友重新在丘比的系统内部行动，漫画便让和美的叙述直白断言，她们不会变成魔女。
《魔法少女小圆》中的希望—绝望循环被彻底否定；
取而代之的，是这样一种对女性力量的描绘：
只有当女性力量被限制在对抗女性力量之负面表达的角色之中时，它才被允许继续存在。
只要魔法少女把自己限制在与魔女战斗之中，不去试图挑战丘比或其腐败的系统，漫画便断言，她们可以继续活下去，并继续运用自己的力量。

归根结底，《魔法少女和美》整齐地落入了那个陷阱：
它（希望只是无意地）断言，丘比的系统是必要的，因为必须用它来限制女性，防止她们以可能挑战现状的方式运用力量。
